	\begin{UseCase}{CU11.1.1}{Gestionar Acciones}{
		
		El colaborador (\hyperlink{jefe}{Líder de Análisis} o \hyperlink{analista}{Analista}) podrá gestionar las \hyperlink{EntidadAccion}{acciones} propias de cada {pantalla}{Pantalla}.\\
		
		Las acciones son componentes como botones, hipervínculos y opciones del menú que cumplen una funcionalidad dentro de la pantalla. Este caso de uso permite al colaborador visualizar en una tabla, el registro de las acciones que hay en la pantalla sobre la cual se está operando, así como solicitar el registro de una acción nueva, modificar alguna existente (teniendo la posibilidad de actualizar algún dato de registro), o eliminarla (en caso de que la acción por algún motivo ya no forme parte de la pantalla)\\
			
     	Las acciones disponibles para cada \hyperlink{EntidadAccion}{acción} dependerán del estado en el que se encuentre el caso de uso donde son referenciadas. 

	}
	\UCitem{Actor}{\hyperlink{jefe}{Líder de análisis}, \hyperlink{analista}{Analista}}
	\UCitem{Propósito}{Proporcionar al actor un mecanismo para llevar el control de las acciones pertenecientes a una pantalla.}
	\UCitem{Entradas}{Ninguna}
	\UCitem{Salidas}{\begin{itemize}
			\item \cdtRef{proyectoEntidad:claveProyecto}{Clave del proyecto}: Lo obtiene el sistema.
			\item \cdtRef{proyectoEntidad:nombreProyecto}{Nombre del proyecto}: Lo obtiene el sistema.
			\item \cdtRef{EntidadPantalla:nombreIU}{Nombre de la pantalla:} Lo obtiene el sistema.
			\item \cdtRef{EntidadAccion}{Acción}: Tabla que muestra \cdtRef{EntidadAccion:imagenACC}{Imagen} y \cdtRef{EntidadAccion:nombreACC}{Número} de todas las acciones registradas en un pantalla.
			\item \cdtIdRef{MSG2}{No existe información}: Se muestra en la pantalla \IUref{IU7.1.1.1}{Gestionar Acciones} cuando no existen acciones registradas.
	\end{itemize}}
	
	\UCitem{Precondiciones}{Que exista al menos una pantalla registrada.}
	\UCitem{Postcondiciones}{Ninguna}
	\UCitem{Errores}{Ninguno}
	\UCitem{Tipo}{Secundario, extiende del caso de uso \UCref{CU11}{Gestionar Pantallas}.}
\end{UseCase}
%--------------------------------------
\begin{UCtrayectoria}
	\UCpaso[\UCactor] Solicita gestionar las acciones de una pantalla seleccionando el icono \raisebox{-1mm}{\includegraphics[height=11pt]{images/Iconos/acciones}} de la pantalla \IUref{IU7}{Gestionar Pantallas}.
	\UCpaso[\UCsist] Obtiene la información de las acciones registradas de la pantalla seleccionada. \hyperlink{CU11-1-1:TAA}{[Trayectoria A]}
	\UCpaso[\UCsist] Ordena las acciones alfabéticamente basándose en el nombres de los mismas.
	\UCpaso[\UCsist] Muestra la información de las acciones en la pantalla \IUref{IU7.1.1.1}{Gestionar Acciones} y las operaciones disponibles de acuerdo a la regla de negocio \BRref{RN15}{Operaciones disponibles}. \label{CU11-1-1-P4}
	\UCpaso[\UCactor] Gestiona las acciones a través de los botones: \IUbutton{Registrar}, \editar y \eliminar. 
\end{UCtrayectoria}		
%--------------------------------------
\hypertarget{CU11-1-1:TAA}{\textbf{Trayectoria alternativa A}}\\
\noindent \textbf{Condición:} No existen registros de acciones.
\begin{enumerate}
	\UCpaso[\UCsist] Muestra el mensaje \cdtIdRef{MSG2}{No existe información} en la pantalla \IUref{IU17A.1.1.1}{Gestionar Acciones: Sin registros} para indicar que no hay registros de acciones para mostrar.  \label{CU11-1-1-TA1}
	\UCpaso[\UCactor] Gestiona las acciones a través del botón: \IUbutton{Registrar}. 
	\item[- -] - - {\em {Fin de la trayectoria}}.%
\end{enumerate}
%--------------------------------------
\subsubsection{Puntos de extensión}

\UCExtenssionPoint{El actor requiere registrar una acción}{Presionando el botón \IUbutton{Registrar} del paso \ref{CU11-1-1-P4} de la trayectoria principal o del paso \ref{CU11-1-1-TA1} de la Trayectoria alternativa A.}{\UCref{CU11.1.1.1}{Registrar Acción}}
\UCExtenssionPoint{El actor requiere modificar una acción}{Presionando el icono \editar del paso \ref{CU11-1-1-P4} de la trayectoria principal.}{\UCref{CU11.1.1.2}{Modificar Acción}}
\UCExtenssionPoint{El actor requiere eliminar una acción}{Presionando el icono \eliminar del paso \ref{CU11-1-1-P4} de la trayectoria principal.}{\UCref{CU11.1.1.3}{Eliminar Acción}}
